\BourbakiExerciseGroup

\begin{enumerate}
\def\labelenumi{\arabic{enumi})}
\tightlist
\item
  (Algebra of sequences of exponential type) Given two integers \(p, q \in \mathbf{N}\) , we denote by \(((p, q))\) the binomial coefficient \(\frac{(p + q)!}{p!q!}\) . Let \(k\) be a commutative ring and \(B\) an associative commutative unital \(k\) -algebra. By a sequence of exponential type of elements of \(B\) we understand a sequence \(a = (a_p)_{p \in \mathbf{N}}\) such that \(a_0 = 1_B\) and \(a_p a_q = ((p, q)) a_{p + q}\) for any \(p, q \in \mathbf{N}\) . The set of all sequences of exponential type of \(B\) is written \(\mathcal{E}(B)\) . \(a)\) Verify that there is on \(\mathcal{E}(B)\) a unique \(k\) -algebra structure such that
\end{enumerate}

\[
\begin{array}{c} (a + b) _ {p} = \sum_ {r + s = p} a _ {r} b _ {s}, \\ (a b) _ {p} = p! a _ {p} b _ {p}, \\ (\lambda a) _ {p} = \lambda^ {p} a _ {p}, \end{array}
\]

for any sequence \(a, b \in \mathcal{E}(\mathbf{B})\) and \(p \in \mathbf{N}, \lambda \in k\) . The mapping \(a \mapsto a_1\) is a homomorphism of the algebra \(\mathcal{E}(\mathbf{B})\) into \(\mathbf{B}\) .

\begin{enumerate}
\def\labelenumi{\alph{enumi})}
\setcounter{enumi}{1}
\item
  Suppose that \(k\) contains a subring isomorphic to \(\mathbf{Q}\) . Show that for every \(x \in B\) the sequence \(f(x)\) defined by \(f(x)_p = \frac{1}{p!} x^p\) is of exponential type. Show that \(f\) is an isomorphism of \(B\) onto \(\mathcal{E}(B)\) .
\item
  Let \(\mathbf{B}'\) be an associative commutative unital \(k\) -algebra and let \(h\) be a homomorphism of \(\mathbf{B}\) into \(\mathbf{B}'\) . For each \(a \in \mathcal{E}(\mathbf{B})\) the sequence \(a'\) defined by \(a_p' = h(a_p)\) for any \(p \in \mathbf{N}\) belongs to \(\mathcal{E}(\mathbf{B}')\) . Show that the mapping \(\mathcal{E}(h): a \mapsto a'\) is a homomorphism of \(\mathcal{E}(\mathbf{B})\) into \(\mathcal{E}(\mathbf{B}')\) . Verify that \(\mathcal{E}(g \circ h) = \mathcal{E}(g) \circ \mathcal{E}(h)\) when \(g\) and \(h\) are composable \(k\) -algebra homomorphisms.
\item
  Let E be a \(k\) -module and \(\mathbf{TS}(\mathbf{E})\) the algebra of symmetric tensors on E. Show that for each \(x \in \mathbf{E}\) , \((\gamma_p(x))\) is a sequence of exponential type in the algebra \(\mathbf{TS}(\mathbf{E})\) (IV, p.~45).
\end{enumerate}

\begin{enumerate}
\def\labelenumi{\arabic{enumi})}
\setcounter{enumi}{1}
\tightlist
\item
  (Functor \(\Gamma\) ) Denote by \(k\) a commutative ring and by E a \(k\) -module. By the gamma algebra of E we understand the associative commutative unital algebra defined by the generating system \(\mathbf{N} \times \mathbf{E}\) and by the relators (cf.~III, p.~450)
\end{enumerate}

\[
\begin{array}{c} (0, x) - 1, \\ (p, \lambda x) - \lambda^ {p} (p, x), \\ (p, x + y) - \sum_ {r + s = p} (r, x) (s, y), \\ (p, x) (q, x) - ((p, q)) (p + q, x), \end{array}
\]

where \(p, q\) range over \(\mathbf{N}, x, y\) range over \(\mathbf{E}\) and \(\lambda\) over \(k\) . This algebra is denoted by \(\Gamma(\mathbf{E})\) . For each \(p \in \mathbf{N}\) we denote by \(\gamma_p\) the mapping of \(\mathbf{E}\) into \(\Gamma(\mathbf{E})\) composed of the injection \(x \mapsto (p, x)\) and the canonical homomorphism of the free commutative algebra on \(\mathbf{N} \times \mathbf{E}\) into \(\Gamma(\mathbf{E})\) .

\begin{enumerate}
\def\labelenumi{\alph{enumi})}
\item
  Verify that for each \(x \in \mathbb{E}\) , the sequence \(\gamma_p(x)\) is a sequence of exponential type of elements of \(\Gamma(\mathbb{E})\) (Exercise 1) and that the mapping \(\gamma : x \mapsto (\gamma_p(x))_{p \in \mathbb{N}}\) is a homomorphism of the \(k\) -module \(\mathbb{E}\) into the \(k\) -algebra \(\mathcal{E}(\Gamma(\mathbb{E}))\) .
\item
  (Universal property of \(\Gamma(\mathrm{E})\) ). Let B be an associative commutative unital k-algebra and \(\varphi\) a homomorphism of the k-module E into \(\mathcal{E}(\mathrm{B})\) . Show that there exists a unique algebra homomorphism \(h: \Gamma(\mathrm{E}) \to \mathrm{B}\) such that \(\varphi = \mathcal{E}(h) \circ \gamma\) .
\item
  Show that there exists a unique unital homomorphism of algebras \(\varepsilon : \Gamma(E) \to k\) such that \(\varepsilon(\gamma_p(x)) = 0\) for all \(p > 0\) and all \(x \in E\) .
\end{enumerate}

\(d)\) For each \(p \in \mathbb{N}\) let \(\Gamma_p(E)\) be the submodule of \(\Gamma(E)\) generated by the elements

\[
\gamma_ {\nu_ {1}} (x _ {1}) \dots \gamma_ {\nu_ {s}} (x _ {s})
\]

where \(s \in \mathbf{N}\) , \(\nu \in \mathbf{N}^s\) and \(|\nu| = \sum_{i} \nu_i = p\) . Verify that the submodules \(\Gamma_p(\mathrm{E})\) form a

graduation of the algebra \(\Gamma(E)\) . Show that \(\gamma_{1}\) is an isomorphism of \(E\) onto \(\Gamma_{1}(E)\) (to verify that \(\gamma_{1}\) is injective, consider the algebra \(B = k \times E\) with the product defined by the formula \((\lambda, x)(\mu, y) = (\lambda \mu, \lambda y + \mu x)\) and show that there exists a homomorphism of \(\Gamma(E)\) into \(B\) which maps \(\gamma_{1}(x)\) to \((0, x)\) for any \(x \in E\) ).

\begin{enumerate}
\def\labelenumi{\alph{enumi})}
\setcounter{enumi}{4}
\item
  Let E and F be two k-modules and h a homomorphism of E into F. Show that there exists a unique homomorphism of graded algebras \(\Gamma(h):\Gamma(\mathrm{E})\to\Gamma(\mathrm{F})\) such that \(\gamma_{p}\circ h=\Gamma(h)\circ\gamma_{p}\) for all \(p\in N\) . Verify that if g and h are homomorphisms of k-modules such that \(g\circ h\) is defined, then \(\Gamma(g\circ h)=\Gamma(g)\circ\Gamma(h)\) .
\item
  Let \(h: \mathbf{E} \to \mathbf{F}\) be a surjective homomorphism of \(k\) -modules. Show that \(\Gamma(h)\) is surjective and that its kernel is the ideal of \(\Gamma(\mathbf{E})\) generated by the elements \(\gamma_p(x)\) with \(p > 0\) and \(x \in \operatorname{Ker}(h)\) .
\item
  Let E and F be two k-modules. Show that there exists a unique homomorphism \(\varphi\) of the algebra \(\Gamma(E \times F)\) into the algebra \(\Gamma(E) \otimes_{k} \Gamma(F)\) such that
\end{enumerate}

\[
(\varphi \circ \gamma_ {p}) (x, y) = \sum_ {r + s = p} \gamma_ {r} (x) \otimes \gamma_ {s} (y)
\]

for any \(p \in \mathbf{N}\) , \(x \in E\) and \(y \in F\) . Show that this homomorphism is an isomorphism compatible with the graduations.

\begin{enumerate}
\def\labelenumi{\arabic{enumi})}
\setcounter{enumi}{2}
\tightlist
\item
  Let \(k\) be a commutative ring and \(E\) a \(k\) -module. Show that there exists a unique homomorphism \(\Delta\) of the algebra \(\Gamma(E)\) (Exercise 2) into the algebra \(\Gamma(E) \otimes_k \Gamma(E)\) such that
\end{enumerate}

\[
(\Delta \circ \gamma_ {p}) (x) = \sum_ {r + s = p} \gamma_ {r} (x) \otimes \gamma_ {s} (x)
\]

for all \(x \in E\) and all \(p \in N\) . Show that the coproduct \(\Delta\) is coassociative, cocommutative (III, p.~580-581) and defines on \(\Gamma(E)\) a graded bigebra structure (III, p.~585). Show that the homomorphism \(\varepsilon\) defined in Exercise 2, c) is a counit.

\begin{enumerate}
\def\labelenumi{\arabic{enumi})}
\setcounter{enumi}{3}
\item
  Let \(E_{1}\) and \(E_{2}\) be two modules over a commutative ring k and let h be a homomorphism of \(E_{1}\) into \(E_{2}\) . Verify that \(\Gamma(h)\) (Exercise 2, e)) is a homomorphism of graded bigebras (cf.~Exercise 3).
\item
  Let \(\Gamma(\mathbf{Z})\) be the gamma algebra of the \(\mathbf{Z}\) -module \(\mathbf{Z}\) (Exercise 2). For each \(p \in \mathbf{N}\) put \(\mathrm{T}^{(p)} = \gamma_p(1)\) .
\end{enumerate}

\begin{enumerate}
\def\labelenumi{\alph{enumi})}
\tightlist
\item
  Show that \((\mathbf{T}^{(p)})_{p\in \mathbb{N}}\) is a basis of the \(\mathbf{Z}\) -module \(\Gamma (\mathbf{Z})\) . Show that there exists a unique homomorphism \(\theta\) of the algebra \(\Gamma (\mathbf{Z})\) into the algebra \(\Gamma (\mathbf{Z})\otimes \Gamma (\mathbf{Z})\) such that
\end{enumerate}

\[
\theta \left(\mathrm{T} ^ {(p)}\right) = p! \mathrm{T} ^ {(p)} \otimes \mathrm{T} ^ {(p)}
\]

for all \(p \in \mathbf{N}\) .

\begin{enumerate}
\def\labelenumi{\alph{enumi})}
\setcounter{enumi}{1}
\item
  Let B be a commutative ring and H the set of homomorphisms of the ring \(\Gamma(\mathbf{Z})\) into the ring B. Show that for every \(f \in H\) the sequence \(f(\mathrm{T}^{(p)})\) , \(p \in N\) is a sequence of exponential type of elements of B. Show that the mapping \(\varphi: f \mapsto (f(\mathrm{T}^{(p)}))_{p \in \mathbb{N}}\) is a bijection of H onto \(\mathcal{E}(\mathbf{B})\) (Exercise 1).
\item
  For \(f, g \in \mathbf{H}\) denote by \(f * g\) the homomorphism of \(\Gamma(\mathbf{Z}) \otimes \Gamma(\mathbf{Z})\) into \(\mathbf{B}\) defined by \((f * g)(u \otimes v) = f(u)g(v)\) for any \(u, v \in \Gamma(\mathbf{Z})\) . Verify that for any \(f, g \in \mathbf{H}\) we have
\end{enumerate}

\[
\begin{array}{c} \varphi (f) + \varphi (g) = \varphi ((f * g) \circ \Delta), \\ \varphi (f) \varphi (g) = \varphi ((f * g) \circ \theta). \end{array}
\]

\begin{enumerate}
\def\labelenumi{\arabic{enumi})}
\setcounter{enumi}{5}
\item
  Let \(n\) be an integer \(> 0\) and let \(\Gamma(\mathbf{Z}/n\mathbf{Z})\) be the gamma algebra of the \(\mathbf{Z}\) -module \(\mathbf{Z}/n\mathbf{Z}\) (Exercise 2). Show that for every \(p > 0\) , \(\Gamma_p(\mathbf{Z}/n\mathbf{Z})\) is isomorphic to the quotient of \(\mathbf{Z}\) by the subgroup generated by the integers \(n^k((k, p - k))\) , where \(1 \leqslant k \leqslant p\) . Show that if \(p\) and \(n\) are distinct primes, then \(\Gamma_p(\mathbf{Z}/n\mathbf{Z})\) is cyclic of order \(n\) . Show that if \(n\) is even, then \(\Gamma_2(\mathbf{Z}/n\mathbf{Z})\) is cyclic of order \(2n\) .
\item
  (Extension of scalars) Let \(k\) be a commutative ring, \(L\) an associative commutative unital \(k\) -algebra and \(E\) a \(k\) -module.
\end{enumerate}

\begin{enumerate}
\def\labelenumi{\alph{enumi})}
\tightlist
\item
  Let \(\Gamma(L \otimes_k E)\) be the gamma algebra of the L-module \(L \otimes_k E\) (Exercise 2). Show that there exists a unique homomorphism \(\theta\) of \(L \otimes_k \Gamma(E)\) into \(\Gamma(L \otimes_k E)\) compatible with the algebra structure on \(L\) and such that
\end{enumerate}

\[
\theta (1 \otimes \gamma_ {p} (x)) = \gamma_ {p} (1 \otimes x)
\]

for all \(x \in \mathbf{E}\) and \(p \in \mathbf{N}\) .

\begin{enumerate}
\def\labelenumi{\alph{enumi})}
\setcounter{enumi}{1}
\tightlist
\item
  Show that \(\theta\) is an isomorphism of graded bigebras over \(L\) .
\end{enumerate}

\begin{enumerate}
\def\labelenumi{\arabic{enumi})}
\setcounter{enumi}{7}
\tightlist
\item
  Let \(k\) be a commutative ring and \(E\) a \(k\) -module.
\end{enumerate}

\begin{enumerate}
\def\labelenumi{\alph{enumi})}
\tightlist
\item
  Show that there exists a unique homomorphism g of the algebra \(\Gamma(\mathrm{E})\) (Exercise 2) into the algebra TS(E) of symmetric tensors over E (IV, p.~43) such that
\end{enumerate}

\[
(g \circ \gamma_ {p}) (x) = x \otimes x \otimes \dots \otimes x \quad (p \text {   factors   } x),
\]

for any \(x \in \mathbb{E}\) , \(p \in \mathbb{N}\) . Show that if \(\mathbb{E}\) is a free \(k\) -module with basis \((e_i)_{i \in I}\) , then \(g\) is an isomorphism of graded bigebras (cf.~Exercise 3 and IV, p.~51) and the elements \(\gamma_\alpha = \prod_i \gamma_{\alpha(i)}(e_i)\) , where \(\alpha \in \mathbf{N}^{(I)}\) , form a basis of \(\Gamma(\mathbb{E})\) .

\begin{enumerate}
\def\labelenumi{\alph{enumi})}
\setcounter{enumi}{1}
\tightlist
\item
  Let \(\xi\) be a mapping of \((1, n)\) into \(\mathbf{N}\) and let \(p = \sum_{i=1}^{n} \xi(i)\) . Show that for every sequence
\end{enumerate}

\(x_{1}, x_{2}, \ldots, x_{n}\) of elements of E, the image of \(\gamma_{\xi(1)}(x_{1}) \ldots \gamma_{\xi(n)}(x_{n})\) under \(g\) is \(\sum_{\varphi} x_{\varphi(1)} \otimes \ldots \otimes x_{\varphi(p)}\) where \(\varphi\) runs over the set of all mappings of \([1, p]\) into \([1, n]\) such that

Card \(\varphi^{-1}(i) = \xi(i)\) for each \(i \in [1, n]\) (cf.~IV, p.~45, Prop. 3).

\begin{enumerate}
\def\labelenumi{\alph{enumi})}
\setcounter{enumi}{2}
\tightlist
\item
  Let \(f\) be the mapping of \(\otimes^p \mathbf{E}\) into \(\Gamma_p(\mathbf{E})\) defined by \(f(x_1 \otimes \ldots \otimes x_p) = \gamma_1(x_1) \ldots \gamma_1(x_p)\) for any \(x_1, \ldots, x_p \in \mathbf{E}\) . Show that for each tensor \(t \in \otimes^p \mathbf{E}\) , \((g \circ f)(t)\) is the symmetrization of \(t\) . Show that for each \(u \in \Gamma_p(\mathbf{E})\) we have \((f \circ g)(u) = p! u\) .
\end{enumerate}

\begin{enumerate}
\def\labelenumi{\arabic{enumi})}
\setcounter{enumi}{8}
\tightlist
\item
  (Polynomial laws) Denote by E and F two modules over the commutative ring k. By a polynomial law of E in F we understand, for each (associative commutative and unital) k-algebra L, a mapping \(\varphi_{L}\) of \(L \otimes_{k} E\) into \(L \otimes_{k} F\) satisfying the condition:
\end{enumerate}

If L and R are two k-algebras and f is a unital homomorphism of L into R, then

\[
(f \otimes 1) \circ \varphi_ {L} = \varphi_ {R} \circ (f \otimes 1).
\]

\begin{enumerate}
\def\labelenumi{\alph{enumi})}
\tightlist
\item
  Let \(\varphi\) be a polynomial law of \(E\) into \(F\) , let \(L = k[(T_i)_{i \in I}]\) be the polynomial algebra in a family of indeterminates \(T_i\) over \(k\) and let \((x_i)_{i \in I}\) be a family of finite support of elements of \(E\) . Show that there exists a unique family of finite support \((y_\xi)_{\xi \in N^{(I)}}\) of elements of \(F\) such that
\end{enumerate}

\[
\varphi_ {L} \left(\sum_ {i} T _ {i} \otimes x _ {i}\right) = \sum_ {\xi \in N ^ {(I)}} T ^ {\xi} \otimes y _ {\xi}.
\]

Show that if \(\mathbf{R}\) is an algebra over \(k\) , then for every family \((a_i)_{i \in I}\) of elements of \(\mathbf{R}\) we have

\[
\varphi_ {R} \left(\sum_ {i} a _ {i} \otimes x _ {i}\right) = \sum_ {\xi \in \mathbf {N} ^ {(I)}} a ^ {\xi} \otimes y _ {\xi}.\tag{1}
\]

\begin{enumerate}
\def\labelenumi{\alph{enumi})}
\setcounter{enumi}{1}
\tightlist
\item
  A family \((\varphi^j)_{j \in \mathbb{J}}\) of polynomial laws of E into F is called summable if for every algebra L over \(k\) and every \(u \in \mathrm{L} \otimes \mathrm{E}\) , the family \(\varphi_L^j(u)\) has finite support. Show that if \((\varphi^j)_{j \in \mathbb{J}}\) is a summable family of polynomial laws of E into F, then there exists a unique polynomial law \(\varphi\) of E into F such that for every \(k\) -algebra L and every \(u \in \mathrm{L} \otimes \mathrm{E}\) we have \(\sum_j \varphi_L^j(u) = \varphi_L(u)\) . This polynomial law is called the sum of the laws \(\varphi^j\) and is denoted by
\end{enumerate}

\[
\sum_ {j} \varphi^ {j}.
\]

\begin{enumerate}
\def\labelenumi{\alph{enumi})}
\setcounter{enumi}{2}
\item
  Let E, F, G be three k-modules, \(\varphi\) a polynomial law of E into F and \(\psi\) a polynomial law of F into G. Show that there exists a unique polynomial law \(\eta\) of E into G such that for every k-algebra L we have \(\eta_{L} = \psi_{L} \circ \varphi_{L}\) . This law \(\eta\) is called the composition of \(\varphi\) and \(\psi\) and is denoted by \(\psi \circ \varphi\) .
\item
  Let \(p\) be an integer \(\geqslant 0\) . A polynomial law \(\varphi\) of \(E\) into \(F\) is said to be homogeneous of degree \(p\) if for every algebra \(L\) over \(k\) we have \(\varphi_{L}(au) = a^{p}\varphi_{L}(u)\) for any \(a \in L\) and \(u \in L \otimes E\) . Show that if \(\varphi\) is homogeneous of degree \(p\) , then in formula (1) we have \(y_{\xi} = 0\) for \(|\xi| = \sum_{i} \xi(i) \neq p\) . Determine all homogeneous polynomial laws of degrees 0 and 1.
\item
  With the same notation as in c), show that if \(\varphi\) and \(\psi\) are homogeneous laws of degrees \(p\) and \(q\) respectively, then \(\psi \circ \varphi\) is a homogeneous law of degree \(pq\) .
\item
  Let \(p\) be an integer \(\geqslant 0\) . For every commutative \(k\) -algebra \(L\) we denote by \(\varepsilon_L\) the injection \(a \mapsto aT\) of \(L\) into the polynomial algebra \(L[T]\) and by \(\beta_L^p\) the homomorphism of the \(L\) -module \(L[T]\) into \(L\) which associates with each element \(\sum_k a_k T^k \in L[T]\) the coefficient \(a_p\) . Let \(\varphi\) be a polynomial law of \(E\) into \(F\) . Show that the
\end{enumerate}

mappings \(\varphi_{L}^{p} = (\beta_{L}^{p}\otimes \mathrm{Id}_{F})\circ \varphi_{L[T]}\circ (\varepsilon_{L}\otimes \mathrm{Id}_{E})\) form a polynomial law of E into F which is homogeneous of degree \(p\) . This polynomial law is called the homogeneous component of degree \(p\) of the law \(\varphi\) . Show that the homogeneous components of \(\varphi\) form a summable family with sum \(\varphi\) .

\begin{enumerate}
\def\labelenumi{\arabic{enumi})}
\setcounter{enumi}{9}
\tightlist
\item
  Denote by \(k\) a commutative ring and by \(E\) a \(k\) -module. For each integer \(p \geqslant 0\) and every \(k\) -algebra \(L\) we denote by \(\gamma_{p,L}\) the mapping of \(L \otimes E\) into \(L \otimes \Gamma_p(E)\) composed of \(\gamma_p: L \otimes E \to \Gamma_p(L \otimes E)\) and the canonical isomorphism of \(\Gamma_p(L \otimes E)\) onto \(L \otimes \Gamma_p(E)\) (cf.~Exercises 2 and 7).
\end{enumerate}

\begin{enumerate}
\def\labelenumi{\alph{enumi})}
\tightlist
\item
  Verify that the mappings \(\gamma_{p,L}\) form a polynomial law (Exercise 9) of E into \(\Gamma_p(E)\) , which is homogeneous of degree \(p\) ; this law will in what follows be denoted by \(\boldsymbol{\gamma}_p\) . Show that for every \(k\) -module F and every \(\sigma \in \operatorname{Hom}(\Gamma_p(E), F)\) the mappings \(\varphi_L = (\operatorname{Id}_L \otimes \sigma) \circ \boldsymbol{\gamma}_{p,L}\) form a polynomial law of E into F which is homogeneous of degree \(p\) . This law is called the law associated with \(\sigma\) .
\end{enumerate}

Let \((x_{i})_{i\in I}\) be a family of elements of E and \((a_{i})_{i\in I}\) a family with finite support of elements of L. Show that

\[
\varphi_ {L} \left(\sum_ {i} a _ {i} \otimes x _ {i}\right) = \sum_ {\xi \in N ^ {(I)}, | \xi | = p} a ^ {\xi} \otimes \sigma (\gamma_ {\xi} (x))
\]

where \(|\xi| = \sum_{i} \xi(i), a^{\xi} = \prod_{i} a_{i}^{\xi(i)}\) and \(\gamma_{\xi}(x) = \prod_{i} \gamma_{\xi(i)}(x_{i})\) .

\begin{enumerate}
\def\labelenumi{\alph{enumi})}
\setcounter{enumi}{1}
\tightlist
\item
  Let \(\varphi\) be a polynomial law of E into a k-module F, homogeneous of degree p.~Show that there exists a unique linear mapping \(\sigma\) of \(\Gamma_{p}(\mathrm{E})\) into F such that \(\varphi\) is the law associated with \(\sigma\) . (Suppose first that E is free on a basis \((x_i)_{i \in I}\) ; using Exercise 9, show that in this case there exists a unique mapping \(\sigma \in \operatorname{Hom}(\Gamma_p(E), F)\) such that Relation (1) of Exercise 9 holds for every algebra L and every family \((a_i)_{i \in I}\) with finite support of elements of L. Treat the general case by introducing a free \(k\) -module \(E'\) and a surjective homomorphism \(h: E' \to E\) . The polynomial law \(\varphi\) , composed with the law defined by \(h\) is a polynomial law of \(E'\) into F which is associated to a linear mapping \(\sigma': \Gamma_p(E') \to F\) ; show that \(\sigma'\) is the composite of \(\Gamma_p(h)\) and a mapping \(\sigma: \Gamma_p(E) \to F\) such that \(\varphi\) is associated with \(\sigma\) (cf.~Exercise 2, \(f\) ).)
\end{enumerate}

\begin{enumerate}
\def\labelenumi{\arabic{enumi})}
\setcounter{enumi}{10}
\tightlist
\item
  We keep the notation of Exercise 10. A linear mapping of \(\Gamma(E)\) into \(F\) will be called a series over \(E\) with values in \(F\) . The series over \(E\) with values in \(F\) form a \(k\) -module denoted by \(S(E, F)\) . For \(n \in N\) we say that the series \(\sigma \in S(E, F)\) is of order \(\geqslant n\) if \(\operatorname{Ker}(\sigma) \supset \Gamma_p(E)\) for all \(p < n\) . In what follows we shall equip \(S(E, F)\) with the topological group structure having as neighbourhood base of 0 the \(k\) -modules consisting of the series of order \(\geqslant n\) ( \(n \in N\) ).
\end{enumerate}

Let \(\varphi\) be a polynomial law of \(E\) into \(F\) . Show that there exists a unique series \(\sigma \in S(E, F)\) such that:

\begin{enumerate}
\def\labelenumi{(\roman{enumi})}
\tightlist
\item
  for every \(k\) -algebra \(\mathbf{L}\) and every \(u \in \mathbf{L} \otimes \mathbf{E}\) the family \((\mathrm{Id}_{\mathbf{L}} \otimes \sigma) \gamma_{p,\mathbf{L}}(u) (p \in \mathbf{N})\) has finite support;
\end{enumerate}

\[
\varphi_ {L} (u) = \sum_ {p} \left(\operatorname{Id} _ {L} \otimes \sigma\right) \circ \gamma_ {p, L} (u). \tag {ii}
\]

(Apply the result of Exercise 10, b)).

This series \(\sigma\) is said to be associated to the law \(\varphi\) .

\begin{enumerate}
\def\labelenumi{\arabic{enumi})}
\setcounter{enumi}{11}
\tightlist
\item
  Let \(k\) be a commutative ring, \(E\) a \(k\) -module and \(B\) a \(k\) -algebra (not necessarily associative). Let \(\Delta\) be the coproduct of \(\Gamma(E)\) (cf.~Exercise 3). For any series \(\sigma, \sigma' \in S(E, B)\) we denote by \(\sigma \cdot \sigma'\) the series \(m \circ (\sigma \otimes \sigma') \circ \Delta\) , where \(m\) denotes the linear mapping of \(B \otimes_k B\) into \(B\) defined by the product in \(B\) .
\end{enumerate}

\begin{enumerate}
\def\labelenumi{\alph{enumi})}
\item
  Show that the mapping \((\sigma, \sigma') \mapsto \sigma \cdot \sigma'\) defines on S(E, B) a \(k\) -algebra structure (cf.~Exercise 11). Show that if B is associative (resp. commutative, unital) then S(E, B) is associative (resp. commutative, unital) (cf.~III, p.~578-584).
\item
  Assume that B is associative and unital. Show that the invertible elements of S(E, B) are the series σ such that σ(1) is invertible in B.
\item
  Suppose that E is a free \(k\) -module and \((e_i)_{i \in I}\) is a basis of E. For each \(\alpha \in \mathbf{N}^{(I)}\) put \(e^{(\alpha)} = \Pi \gamma_{\alpha(i)}(e_i)\) and denote by \(f^\alpha\) the series over E with values in \(k\) defined by \(f^\alpha(e^{(\beta)}) = \delta_{\alpha, \beta}\) (Kronecker index) for each \(\beta \in \mathbf{N}^{(I)}\) (cf.~Exercise 8). Show that \(f^\alpha f^\beta = f^{\alpha + \beta}\) for any \(\alpha, \beta \in \mathbf{N}^{(I)}\) . Deduce that when I is finite, the algebra \(S(E, k)\) is isomorphic to the algebra of formal power series \(k[[(X_i)_{i \in I}]\) .
\end{enumerate}

\begin{enumerate}
\def\labelenumi{\arabic{enumi})}
\setcounter{enumi}{12}
\tightlist
\item
  (Divided powers of a series.) The notation is as in Exercises 10 and 11. Let E and F be two \(k\) -modules and \(\sigma\) a series over E with values in F. Denote by \(\sigma\) the homogeneous polynomial law of degree 1 of \(\Gamma(E)\) into F defined by \(\sigma\) . For any integers \(p\) , \(m \in \mathbf{N}\) we denote by \(\sigma_m^{(p)}\) the series associated with the polynomial law \(\gamma_p \circ \sigma \circ (\gamma_0 + \gamma_1 + \ldots + \gamma_m)\) (cf.~Exercise 10).
\end{enumerate}

\begin{enumerate}
\def\labelenumi{\alph{enumi})}
\item
  Show that as \(m\) tends to \(\infty\) , the sequence \(\sigma_{m}^{(p)}\) converges in \(S(E, \Gamma_p(F))\) to a series \(\sigma^{(p)}\) . This series \(\sigma^{(p)}\) is called the divided \(p\) -th power of the series \(\sigma\) .
\item
  Show that if \((x_{i})_{i\in I}\) is a family of elements of \(E\) and \(\eta \in \mathbf{N}^{(I)}\) , then
\end{enumerate}

\[
\sigma^ {(p)} \left(\gamma_ {\eta} (x)\right) = \sum_ {\xi \in \mathscr {L}} \prod_ {\alpha \in N ^ {(I)}} \gamma_ {\xi (\alpha)} \left(\sigma \left(\gamma_ {\alpha} (x)\right)\right),
\]

where \(\mathcal{L}\) is the set of all mappings with finite support of \(\mathbf{N}^{(I)}\) into \(\mathbf{N}\) satisfying the conditions \(\sum_{\alpha} \xi(\alpha) = p\) and \(\sum_{\alpha} \xi(\alpha) \alpha = \eta\) (Apply Ex. 10, \(a\) ).

\begin{enumerate}
\def\labelenumi{\alph{enumi})}
\setcounter{enumi}{2}
\tightlist
\item
  Show that if \(\sigma, \tau \in S(E, F)\) , then for each \(p \in \mathbf{N}\) we have
\end{enumerate}

\[
\left(\sigma + \tau\right) ^ {(p)} = \sum_ {r + s = p} \sigma^ {(r)} \sigma^ {(s)}
\]

(product calculated in the algebra \(S(E, \Gamma(E))\) , cf.~Exercise 12).

\(d\) ) Show that for every integer \(p\in \mathbf{N}\) we have \(\sigma^p = p!\sigma^{(p)}\) , where \(\sigma^p\) is the \(p\) -th power of \(\sigma\) in the algebra \(\mathrm{S}(E, \Gamma(F))\)

\begin{enumerate}
\def\labelenumi{\alph{enumi})}
\setcounter{enumi}{4}
\tightlist
\item
  Let \(\Delta_{\mathrm{E}}\) and \(\Delta_{\mathrm{F}}\) be the respective coproducts in \(\Gamma(\mathrm{E})\) and \(\Gamma(\mathrm{F})\) . Show that for all \(p \in \mathbf{N}\) we have
\end{enumerate}

\[
\Delta_ {F} \circ \sigma^ {(p)} = \sum_ {r + s = p} \left(\sigma^ {(r)} \otimes \sigma^ {(s)}\right) \circ \Delta_ {E}.
\]

\begin{enumerate}
\def\labelenumi{\arabic{enumi})}
\setcounter{enumi}{13}
\tightlist
\item
  With the notation as in Ex. 13, suppose that \(\sigma \in S(E, F)\) is of order \(\geqslant 1\) .
\end{enumerate}

\begin{enumerate}
\def\labelenumi{\alph{enumi})}
\tightlist
\item
  Show that for each \(p \in \mathbf{N}\) , \(\sigma^{(p)}\) is of order \(\geqslant p\) . Deduce that the family \((\sigma^{(p)})_{p \in \mathbf{N}}\) is summable. Writing \(\mathbf{A}(\sigma) = \sum_{p} \sigma^{(p)}\) , verify that for every \(n \in \mathbf{N}\) ,
\end{enumerate}

\[
\mathbf {A} (\sigma) \left(\Gamma_ {n} (\mathrm{E})\right) \subset \Gamma_ {0} (\mathrm{F}) + \dots + \Gamma_ {n} (\mathrm{F}).
\]

\begin{enumerate}
\def\labelenumi{\alph{enumi})}
\setcounter{enumi}{1}
\item
  Show that if \(\operatorname{Ker}(\sigma) \supset \Gamma_m(E)\) for all \(m \neq 1\) , then \(A(\sigma) = \Gamma(\sigma \circ \gamma_1)\) .
\item
  Show that \(\mathbf{A}(\sigma)\) is a homomorphism of the cogebra \(\Gamma(\mathbf{E})\) into the cogebra \(\Gamma(\mathbf{F})\) , compatible with the counits (cf.~Ex. 13, e)).
\item
  Show that if \(\sigma, \tau \in S(E, F)\) are two series of order \(\geqslant 1\) , then \(A(\sigma + \tau) = A(\sigma)A(\tau)\) , the product \(A(\sigma)A(\tau)\) being defined by the product in \(\Gamma(E)\) (cf.~Ex. 13, c)).
\end{enumerate}

\begin{enumerate}
\def\labelenumi{\arabic{enumi})}
\setcounter{enumi}{14}
\tightlist
\item
  With the notation as in Ex. 14, let E, F, G be three k-modules, \(\sigma \in S(E, F)\) , \(\tau \in S(F, G)\) , and suppose that \(\sigma\) is of order \(\geqslant 1\) . The series \(\tau \circ A(\sigma) \in S(E, G)\) is called the composition of \(\sigma\) and \(\tau\) ; it is denoted by \(\tau \circ \sigma\) .
\end{enumerate}

\begin{enumerate}
\def\labelenumi{\alph{enumi})}
\item
  Let \(\varphi\) be a polynomial law of E into F (Ex. 9) and \(\psi\) a polynomial law of F into G, and denote by \(\sigma_{\varphi}\) and \(\sigma_{\psi}\) the series associated with \(\varphi\) and \(\psi\) respectively. Show that if the homogeneous component of degree 0 of \(\varphi\) is zero, then \(\sigma_{\varphi}\) is of order \(\geqslant 1\) and \(\sigma_{\psi} \circ \sigma_{\varphi}\) is the series associated with the polynomial law \(\psi \circ \varphi\) .
\item
  Show that if \(\sigma\) is of order \(\geqslant p\) and \(\tau\) is of order \(\geqslant q\) , then \(\tau \circ \sigma\) is of order \(\geqslant pq\) .
\item
  Show that for every integer \(p\) , \((\tau \circ \sigma)^{(p)} = \tau^{(p)} \circ A(\sigma)\) . Deduce that if \(\tau\) is of order \(\geqslant 1\) and if \(H\) is a \(k\) -module, then for any \(\rho \in S(G, H)\) we have \((\rho \circ \tau) \circ \sigma = \rho \circ (\tau \circ \sigma)\) . d) Suppose that \(G\) is a \(k\) -algebra. Show that the mapping \(\tau \mapsto \tau \circ A(\sigma)\) is a homomorphism of the algebra \(S(F, G)\) into the algebra \(S(E, G)\) .
\end{enumerate}

\begin{enumerate}
\def\labelenumi{\arabic{enumi})}
\setcounter{enumi}{15}
\tightlist
\item
  Let E and F be two modules over a commutative ring \(k\) and let \(p\) be an integer \(\geqslant 0\) .
\end{enumerate}

\begin{enumerate}
\def\labelenumi{\alph{enumi})}
\item
  Show that if \(p \leqslant 2\) , the mapping \(\lambda_p\) of \(\operatorname{Hom}(\Gamma_p(E), F)\) into \(F^E\) which associates with \(\sigma\) the mapping \(\sigma \circ \gamma_p\) is an injective mapping.
\item
  Show that \(\lambda_1\) has as image \(\operatorname{Hom}_k(E, F)\) . Show that \(\lambda_2\) has as image the set of all mappings \(g\) of \(E\) into \(F\) such that:
\end{enumerate}

\begin{enumerate}
\def\labelenumi{(\roman{enumi})}
\item
  \(g(ax) = a^2 g(x)\) for any \(a \in k\) and \(x \in E\) ;
\item
  the mapping \((x, y) \mapsto g(x + y) - g(x) - g(y)\) is a bilinear mapping of \(E \times E\) into \(F\) . c) Show that if \(E\) is a free \(k\) -module, then the image of \(\lambda_p\) is the module \(\operatorname{Pol}_k^p(E, F)\) of homogeneous polynomial mappings of degree \(p\) of \(E\) into \(F\) .
\end{enumerate}

